% chapters/multimodal.tex — Multimodal chapter (narrative text only)
\chapter{Multimodal Agentic Reasoning}
\label{chap:multimodal}

Agentic systems are increasingly required to reason over heterogeneous
modalities — text, images, audio, and structured data — within a single
interaction~\cite{yao2025multimodalsurvey}.
Multimodal \LLM{}s unify these modalities in a shared representation space,
enabling cross-modal evidence fusion.
\index{multimodal reasoning}
\index{cross-modal evidence}

\section{Agentic Multimodal Architectures}
\label{sec:multimodal-arch}

\cite{yao2025multimodalsurvey} survey the design space of agentic multimodal
\LLM{}s along four axes: perception (how modalities are encoded), grounding
(how encodings are aligned to language), reasoning (how evidence from multiple
modalities is fused), and action (how the agent produces outputs or invokes
tools based on multimodal context).

The key architectural choice is whether modalities are fused early (at the
embedding layer), late (at the output layer), or dynamically (the model
selects a fusion strategy per step).
Dynamic fusion consistently outperforms static alternatives on tasks that
require selective attention to a single modality.

\section{Proactive Multimodal Retrieval}
\label{sec:multimodal-retrieval}

As discussed in Chapter~\ref{chap:retrieval}, proactive retrieval strategies
are especially effective in multimodal settings, where the cost of
retrieving an image is substantially higher than retrieving a text
snippet~\cite{wang2025proactiveretrievalmedical}.
The routing model must therefore weight the expected information gain
against the retrieval cost when deciding whether to fetch a visual asset.

\section{Open Challenges}
\label{sec:multimodal-challenges}

Key open problems include: (i)~compositional reasoning across more than two
modalities simultaneously; (ii)~long-context multimodal memory (connecting
Chapter~\ref{chap:memory} to this chapter); and (iii)~evaluation protocols
that test multimodal reasoning independently of language understanding.
